\documentclass[11pt]{article}
\usepackage[margin=2.6cm]{geometry}
\usepackage{amsmath,amssymb}
\usepackage{bm}
\usepackage[T1]{fontenc}
\usepackage[utf8]{inputenc}
\usepackage{parskip}
\usepackage{microtype}
\newcommand{\neff}{n_{\mathrm{eff}}}
\newcommand{\R}[1]{\medskip\noindent\textbf{#1}\quad}
\newenvironment{rquote}{\begin{quote}\itshape}{\end{quote}}
\begin{document}
\begin{center}
{\LARGE\bfseries Response to Reviewers}\\[6pt]
{\large Manuscript MATG-D-26-00180 --- \emph{Mathematical Geosciences},
third revision}\\[4pt]
``Effective-sample-size calibration of a modality test for cluster homogeneity
and the number of classes in spatially autocorrelated data''\\[4pt]
Francisco Parra \,\textperiodcentered\, Universidad de Santiago de Chile
\end{center}
\medskip

I thank the Editor for the decision and Reviewer~1 for a third careful
reading. Every point has been implemented. Reviewer comments are quoted in
italics; section and equation numbers refer to the revised manuscript, and the
changes made in this round are shown in blue in the marked copy.

\section*{Reviewer 1}

\R{1. Maximising a distribution.}
\begin{rquote}
``Sec.~3.4, l.~2/3 `least favourable among unimodal laws in the sense that it
maximises the asymptotic distribution of the rescaled dip' --- you can
maximise a number but not a distribution.''
\end{rquote}
Corrected. The sentence now says that, asymptotically, the rescaled dip of the
uniform law is stochastically largest, so that its $(1-\alpha)$ quantile is at
least that of any other unimodal law in the class considered. It has also been
moved out of the definition of $q_{1-\alpha}(m)$ into a sentence of its own.

\R{2. The limit in Proposition~1.}
\begin{rquote}
``If the limsup in (13) is 0, then write `$=0$' rather than `$\le\alpha$',
which makes the awkward wording below it superfluous.''
\end{rquote}
Done. Equation~(13) now states that the limit is zero, and the sentence that
followed it has been removed. The proof in Section~S1 of the supplement has
been aligned with the new statement.

\R{3. The wording of Condition (C3), and the Proposition itself.}
\begin{rquote}
``In my view `Condition (C3) (\dots) assumes what one would most want proved'
is correct but super awkward as well \dots\ If you feel comfortable with the
awkward state of affairs as stated above, fair enough.''
\end{rquote}
I have kept the Proposition, which is shorter now that point~2 is
implemented, and rewritten the sentence. Remark~1(i) now reads: ``Condition
(C3) carries the analytic weight: it is the empirical-process form of the
effective-sample-size heuristic. We take it as a hypothesis and do not derive
it, so the proposition shows only that the rest of the argument holds once
that heuristic is granted.'' Should the Editor prefer the paper without the
Proposition, it can move entirely to the supplement at the proof stage.

\R{4. The projection direction and the empirical level.}
\begin{rquote}
``p.~16: `\dots\ the second does not inflate the empirical level.' The latter
statement is not proved in general as far as I can see, so don't overclaim.''
\end{rquote}
Corrected. Remark~1(iv) now says that we have no proof for the data-dependent
direction, only the empirical observation that it did not inflate the level in
the experiment of Figure~1a. The same claim appeared in two further places,
Section~3.3 and the discussion of Figure~1 in Section~3.9; both are now
restricted in the same way to the experiment that supports them.

\R{5. English language and bullet-point lists.}
\begin{rquote}
``The manuscript needs further English language improvement (examples below).
The use of bullet-point lists in select sections should also be
restructured.''
\end{rquote}
The three sentences cited have been corrected as indicated: the first is no
longer a fragment (``We recast the choice of $K$ as a recursive homogeneity
test, which asks at each step whether \dots''), and the missing commas have
been added to the other two. The list of contributions in Section~1, the only
bulleted list in the paper, is now a paragraph of prose. Beyond the examples,
the whole manuscript has been revised for language, with attention to the
kinds of problem the examples illustrate: sentence fragments, missing commas
after introductory phrases, dangling modifiers, absolute constructions, and
long sentences that read better as two. These edits do not change the
content, and British spelling is now used consistently.

\end{document}
